% =============================================================================
% Abstract
% Pit-Viper-Inspired Infrared-Thermal and Visual Sensor Fusion
% for Nocturnal UAV Navigation
% =============================================================================
\selectlanguage{hebrew}


מאמר זה מציג מערכת מיזוג תרמי-חזותי חדשנית לניווט רחפנים אוטונומי בתנאי לילה, בהשראת מנגנון קליטת האינפרה-אדום של הצפע הגומתי (\emph{Crotalus cerastes}). המערכת, הקרויה \en{ThermalVisionFusion} (\en{TVF}), משלבת מידע ממצלמה חזותית (\en{RGB}) וממצלמה תרמית (\en{IR}) באמצעות מנגנון קשב צולב אדפטיבי המתאים את משקל המיזוג בהתאם לתנאי התאורה הסביבתיים. האלגוריתם כולל מיפוי חום תרמי כמפת סליאנס ביולוגית לזיהוי מכשולים, תכנון מסלול מבוסס מפת עלות משולבת, ולמידת חיזוק עמוקה להתאמה דינמית לתנאי הסביבה. המערכת מומשה על פלטפורמת \en{NVIDIA Jetson Orin NX} עם צריכת הספק ממוצעת של 22.3 ואט וזמן השהיה כולל של 26.8 מילישניות. תוצאות ניסוי על 12 משימות טיסה ליליות בשטח פתוח מראות שיפור של 73\% בשגיאת המסלול המוחלטת (\en{ATE}) לעומת מערכת \en{ORB-SLAM3} חזותית בלבד (1.15 מ' לעומת 4.28 מ'), שיעור הימנעות ממכשולים של 94.7\%, ושיעור אבדן עקיבה של 3.1\% בלבד. מחקר זה מדגים את הפוטנציאל הרב של מיזוג תרמי-חזותי בהשראת ביולוגית לניווט אוטונומי בתנאי לילה.
